\chapter*{Kết luận}
\addcontentsline{toc}{chapter}{Kết luận}

Việc triển khai CIS Ubuntu Benchmark với Ansible automation đã chứng minh hiệu quả vượt trội so với phương pháp thủ công. Đồ án đã thành công xây dựng một framework hoàn chỉnh cho security hardening, đạt được \textasciitilde 92\% compliance score và tạo foundation cho continuous security improvement.

\section*{Đóng góp chính}
\begin{itemize}
    \item Xây dựng complete automation framework cho CIS compliance
    \item Demonstrated measurable security improvements (+17\% từ baseline)
    \item Created reusable, scalable solution cho production environments
    \item Established baseline cho continuous security monitoring
\end{itemize}

\section*{Giá trị thực tiễn}
\begin{itemize}
    \item \textbf{Risk reduction}: Giảm đáng kể attack surface
    \item \textbf{Compliance achievement}: Đáp ứng industry standards
    \item \textbf{Operational efficiency}: Giảm 95\% manual effort
    \item \textbf{Scalability}: Sẵn sàng cho enterprise deployment
\end{itemize}

\section*{Hướng phát triển}
Framework này có thể mở rộng thành enterprise-grade security automation platform với:
\begin{itemize}
    \item CI/CD security pipeline integration
    \item SIEM connectivity cho advanced threat detection
    \item Auto-remediation capabilities
    \item Continuous compliance monitoring
\end{itemize}

\section*{Final Recommendation}
Organizations nên adopt automation-first approach cho security compliance để ensure consistency, scalability, và continuous improvement trong security posture. Framework này cung cấp foundation vững chắc cho journey towards zero-trust security architecture.

\paragraph{Tầm nhìn tương lai:}
Với sự phát triển của cloud computing và containerization, framework này có thể mở rộng để hỗ trợ multi-cloud environments, Kubernetes security, và DevSecOps practices, tạo thành comprehensive security automation platform cho modern IT infrastructure.
